% depSip.tex -- result usage, with definitions and applications listed separately.
\section{Dependencies between the results}\label{s:dep}
Figure~\ref{fig:dep} records use of a proved conclusion or criterion in
another result, including use without an explicit internal citation.
An arrow $A\to B$ means that the proof of $B$ uses that conclusion of $A$.
Panels group related questions, not connected components of the proof graph;
placing two results together does not assert a proof dependency.

\begin{figure}[H]
\centering
\begin{minipage}{\linewidth}
\centering\textbf{Arbitrary multistability: gain bounds and fixed siphons}\par
{\small (abstract items 1 and 2)}\par\smallskip
\begin{tikzpicture}[x=1cm,y=-1.25cm,font=\footnotesize,
 T/.style={draw,rectangle,inner sep=4pt,fill=blue!6,align=center},
 P/.style={draw,rounded corners=3pt,inner sep=3pt,fill=green!8},
 C/.style={draw,rounded corners=3pt,inner sep=3pt},
 arr/.style={-{Stealth[length=4pt]},thin,gray!70!black}]
\node[C] (cEnt) at (-2.5,2) {\hyperref[c:Ent]{C\,\ref*{c:Ent}}};
\node[P] (pColWie) at (-2.5,0) {\hyperref[p:ColWie]{P\,\ref*{p:ColWie}}};
\node[T] (tKBsharp) at (-2.5,1) {\hyperref[t:KBsharp]{T\,\ref*{t:KBsharp}}\\Sharp uniform NGM bound};
\node[T] (tsiphMulti) at (2.5,1) {\hyperref[t:siphMulti]{T\,\ref*{t:siphMulti}}\\Fixed siphons and transverse matrices};
\draw[arr] (pColWie.south) -- (tKBsharp.north);
\draw[arr] (tKBsharp.south) -- (cEnt.north);
\end{tikzpicture}\par
{\footnotesize Independent constructions; neither theorem uses the other.}
\end{minipage}\par\bigskip
\ifdefined\withSpatial
\begin{minipage}{\linewidth}
\centering\textbf{Diffusion and spreading}\par
{\small (abstract item 3)}\par\smallskip
\begin{tikzpicture}[x=1.2cm,y=-1.5cm,font=\footnotesize,
 T/.style={draw,rectangle,inner sep=3pt,fill=blue!6},
 P/.style={draw,rounded corners=3pt,inner sep=3pt,fill=green!8},
 C/.style={draw,rounded corners=3pt,inner sep=3pt},
 arr/.style={-{Stealth[length=4pt]},thin,gray!70!black}]
\node[C] (cDifSta) at (0,0) {\hyperref[c:DifSta]{C\,\ref*{c:DifSta}}};
\node[T] (tKBsharp) at (1.2,0) {\hyperref[t:KBsharp]{T\,\ref*{t:KBsharp}}};
\node[P] (tspread) at (3,0) {\hyperref[t:spread]{P\,\ref*{t:spread}}};
\node[P] (pOdd) at (4.8,0) {\hyperref[p:Odd]{P\,\ref*{p:Odd}}};
\node[T] (tSpe) at (6,0) {\hyperref[t:Spe]{T\,\ref*{t:Spe}}};
\node[T] (tCmp) at (7,0) {\hyperref[t:Cmp]{T\,\ref*{t:Cmp}}};
\node[T] (tEqWgt) at (8,0) {\hyperref[t:EqWgt]{T\,\ref*{t:EqWgt}}};
\node[P] (pComBas) at (8.9,0) {\hyperref[p:ComBas]{P\,\ref*{p:ComBas}}};
\node[P] (pLVboundary) at (9.9,0) {\hyperref[p:LVboundary]{P\,\ref*{p:LVboundary}}};
\node[T] (tSpdBd) at (0.6,1) {\hyperref[t:SpdBd]{T\,\ref*{t:SpdBd}}};
\node[T] (tSpdQC) at (2.2,1) {\hyperref[t:SpdQC]{T\,\ref*{t:SpdQC}}};
\node[P] (pWeiSpe) at (3.4,1) {\hyperref[p:WeiSpe]{P\,\ref*{p:WeiSpe}}};
\node[C] (cOddSpd) at (4.8,1) {\hyperref[c:OddSpd]{C\,\ref*{c:OddSpd}}};
\node[C] (cDifAll) at (7,1) {\hyperref[c:DifAll]{C\,\ref*{c:DifAll}}};
\node[P] (pDifLya) at (9.4,1) {\hyperref[p:DifLya]{P\,\ref*{p:DifLya}}};
\draw[arr] (cDifSta) -- (tSpdBd);
\draw[arr] (tKBsharp) -- (tSpdBd);
\draw[arr] (tspread) -- (tSpdBd);
\draw[arr] (tspread) -- (tSpdQC);
\draw[arr] (tspread) -- (pWeiSpe);
\draw[arr] (tspread) -- (cOddSpd);
\draw[arr] (pOdd) -- (cOddSpd);
\draw[arr] (tSpe) -- (cDifAll);
\draw[arr] (tCmp) -- (cDifAll);
\draw[arr] (tEqWgt) -- (cDifAll);
\draw[arr] (pComBas) -- (pDifLya);
\draw[arr] (pLVboundary) -- (pDifLya);
\draw[arr] (cDifSta) to[out=25,in=155,looseness=.45] (cDifAll);
\end{tikzpicture}\par
{\footnotesize C\,\ref{c:DifSta} has its own comparison proof and gives item 2 of T\,\ref{t:SpdBd}; P\,\ref{t:spread} gives item 3, and item 1 is proved directly. Its statement applies to the equilibria constructed in T\,\ref{t:siphMulti}, T\,\ref{t:KBsharp}, T\,\ref{t:Net} and T\,\ref{t:KBoverlapCusp}, which are not drawn as arrows. Items of C\,\ref{c:DifAll} use T\,\ref{t:Spe}, T\,\ref{t:Cmp} and T\,\ref{t:EqWgt} in turn; P\,\ref{p:ComBas} gives item 3 of P\,\ref{p:DifLya}, and P\,\ref{p:LVboundary} gives item 4.}
\end{minipage}\par\bigskip
\fi
\begin{minipage}[t]{0.48\linewidth}
\centering\textbf{Reduction tools and their applications}\par
{\small (abstract items \ifdefined\withSpatial 2 and 4\else 2 and 3\fi)}\par\smallskip
\begin{tikzpicture}[x=1.65cm,y=-1.12cm,font=\footnotesize,
 T/.style={draw,rectangle,inner sep=3pt,fill=blue!6},
 P/.style={draw,rounded corners=3pt,inner sep=3pt,fill=green!8},
 C/.style={draw,rounded corners=3pt,inner sep=3pt},
 arr/.style={-{Stealth[length=4pt]},thin,gray!70!black}]
\node[P] (pKBpositive) at (1,0) {\hyperref[p:KBpositive]{P\,\ref*{p:KBpositive}}};
\node[T] (tSpe) at (1,1) {\hyperref[t:Spe]{T\,\ref*{t:Spe}}};
\node[P] (pColWie) at (0,1) {\hyperref[p:ColWie]{P\,\ref*{p:ColWie}}};
\node[P] (pKBshared) at (0,2) {\hyperref[p:KBshared]{P\,\ref*{p:KBshared}}};
\node[T] (tKBhom) at (1,2) {\hyperref[t:KBhom]{T\,\ref*{t:KBhom}}};
\node[T] (tKBexchange) at (2,2) {\hyperref[t:KBexchange]{T\,\ref*{t:KBexchange}}};
\draw[arr] (pKBpositive) -- (tSpe);
\draw[arr] (tSpe) -- (pKBshared);
\draw[arr] (pColWie) -- (pKBshared);
\draw[arr] (tSpe) -- (tKBhom);
\draw[arr] (tSpe) -- (tKBexchange);
\end{tikzpicture}
\end{minipage}\hfill
\begin{minipage}[t]{0.48\linewidth}
\centering\textbf{Cusp criteria and explicit constructions}\par
{\small (abstract items \ifdefined\withSpatial 2, 5, 6\else 2, 4, 5\fi)}\par\smallskip
\begin{tikzpicture}[x=1.65cm,y=-1.12cm,font=\footnotesize,
 T/.style={draw,rectangle,inner sep=3pt,fill=blue!6},
 P/.style={draw,rounded corners=3pt,inner sep=3pt,fill=green!8},
 C/.style={draw,rounded corners=3pt,inner sep=3pt},
 arr/.style={-{Stealth[length=4pt]},thin,gray!70!black}]
\node[P] (psadCus) at (0,0) {\hyperref[p:sadCus]{P\,\ref*{p:sadCus}}};
\node[P] (pKBpositive) at (1,0) {\hyperref[p:KBpositive]{P\,\ref*{p:KBpositive}}};
\node[P] (pQua) at (2,1) {\hyperref[p:Qua]{P\,\ref*{p:Qua}}};
\node[T] (tActF) at (0,1) {\hyperref[t:ActF]{T\,\ref*{t:ActF}}};
\node[T] (tcuspAsym) at (0,2) {\hyperref[t:cuspAsym]{T\,\ref*{t:cuspAsym}}};
\node[T] (tKBhom) at (0,3) {\hyperref[t:KBhom]{T\,\ref*{t:KBhom}}};
\node[T] (tNet) at (1,3) {\hyperref[t:Net]{T\,\ref*{t:Net}}};
\node[T] (tKBoverlapCusp) at (2,3) {\hyperref[t:KBoverlapCusp]{T\,\ref*{t:KBoverlapCusp}}};
\node[C] (cNin) at (2,4) {\hyperref[c:Nin]{C\,\ref*{c:Nin}}};
\draw[arr] (psadCus) -- (tActF);
\draw[arr] (pKBpositive) -- (tActF);
\draw[arr] (psadCus) to[bend right=35] (tcuspAsym);
\draw[arr] (psadCus) to[bend right=35] (tKBhom);
\draw[arr] (psadCus) -- (pQua);
\draw[arr] (psadCus) -- (tNet);
\draw[arr] (pQua) -- (tNet);
\draw[arr] (pQua) to[bend right=35] (tKBoverlapCusp);
\draw[arr] (psadCus) -- (tKBoverlapCusp);
\draw[arr] (tKBoverlapCusp) -- (cNin);
\end{tikzpicture}
\end{minipage}\par\bigskip
\begin{minipage}[t]{0.48\linewidth}
\centering\textbf{Production representation, siphon structure and boundary stability}\par
{\small (abstract items \ifdefined\withSpatial 1, 2, 5, 6\else 1, 2, 4, 5\fi)}\par\smallskip
\begin{tikzpicture}[x=1.65cm,y=-1.12cm,font=\footnotesize,
 T/.style={draw,rectangle,inner sep=3pt,fill=blue!6},
 P/.style={draw,rounded corners=3pt,inner sep=3pt,fill=green!8},
 C/.style={draw,rounded corners=3pt,inner sep=3pt},
 arr/.style={-{Stealth[length=4pt]},thin,gray!70!black}]
\node[P] (pzeroset) at (0,0) {\hyperref[p:zeroset]{P\,\ref*{p:zeroset}}};
\node[P] (pTri) at (1,0) {\hyperref[p:Tri]{P\,\ref*{p:Tri}}};
\node[P] (pKBpositive) at (2,0) {\hyperref[p:KBpositive]{P\,\ref*{p:KBpositive}}};
\node[P] (tCon) at (0,1) {\hyperref[p:toolConservation]{P\,\ref*{p:toolConservation}}};
\node[P] (pInvF) at (1,1) {\hyperref[p:InvF]{P\,\ref*{p:InvF}}};
\node[T] (tActF) at (2,1) {\hyperref[t:ActF]{T\,\ref*{t:ActF}}};
\node[P] (pKBshared) at (2,2) {\hyperref[p:KBshared]{P\,\ref*{p:KBshared}}};
\node[P] (ptraFac) at (0,3) {\hyperref[p:traFac]{P\,\ref*{p:traFac}}};
\node[T] (tGraph) at (1,3) {\hyperref[t:Graph]{T\,\ref*{t:Graph}}};
\node[P] (pOdd) at (2,3) {\hyperref[p:Odd]{P\,\ref*{p:Odd}}};
\node[T] (tNet) at (1,4) {\hyperref[t:Net]{T\,\ref*{t:Net}}};
\node[T] (tKBoverlapCusp) at (2,4) {\hyperref[t:KBoverlapCusp]{T\,\ref*{t:KBoverlapCusp}}};
\node[P] (pKBrepresentation) at (0,4) {\hyperref[p:KBrepresentation]{P\,\ref*{p:KBrepresentation}}};
\draw[arr] (pzeroset) -- (tCon);
\draw[arr] (pTri) -- (pInvF);
\draw[arr] (pTri) -- (ptraFac);
\draw[arr] (pTri) to[bend right=35] (tGraph);
\draw[arr] (pTri) to[bend right=35] (tNet);
\draw[arr] (pTri) -- (tKBoverlapCusp);
\draw[arr] (pKBpositive) -- (tActF);
\draw[arr] (pKBpositive) to[bend right=35] (pKBshared);
\draw[arr] (pInvF) -- (pOdd);
\end{tikzpicture}
\end{minipage}\hfill
\begin{minipage}[t]{0.48\linewidth}
\centering\textbf{Positive weights and contraction}\par
{\small (abstract item \ifdefined\withSpatial 7\else 6\fi)}\par\smallskip
\begin{tikzpicture}[x=1.65cm,y=-1.12cm,font=\footnotesize,
 T/.style={draw,rectangle,inner sep=3pt,fill=blue!6},
 P/.style={draw,rounded corners=3pt,inner sep=3pt,fill=green!8},
 C/.style={draw,rounded corners=3pt,inner sep=3pt},
 arr/.style={-{Stealth[length=4pt]},thin,gray!70!black}]
\node[P] (pColWie) at (0,0) {\hyperref[p:ColWie]{P\,\ref*{p:ColWie}}};
\node[P] (pVarCert) at (2,0) {\hyperref[p:VarCert]{P\,\ref*{p:VarCert}}};
\node[T] (tCmp) at (0,1) {\hyperref[t:Cmp]{T\,\ref*{t:Cmp}}};
\node[P] (pSymEne) at (0,2) {\hyperref[p:SymEne]{P\,\ref*{p:SymEne}}};
\node[P] (pComLya) at (0,3) {\hyperref[p:ComLya]{P\,\ref*{p:ComLya}}};
\node[P] (pComBas) at (1,3) {\hyperref[p:ComBas]{P\,\ref*{p:ComBas}}};
\node[P] (pLVboundary) at (2,3) {\hyperref[p:LVboundary]{P\,\ref*{p:LVboundary}}};
\draw[arr] (pComLya) -- (pComBas);
\draw[arr] (pComBas) -- (pLVboundary);
\node[T] (tWgt) at (2,1) {\hyperref[t:Wgt]{T\,\ref*{t:Wgt}}};
\node[T] (tEqWgt) at (2,2) {\hyperref[t:EqWgt]{T\,\ref*{t:EqWgt}}};
\draw[arr] (pColWie) -- (tCmp);
\draw[arr] (pColWie) -- (tEqWgt);
\draw[arr] (pVarCert) -- (tWgt);
\draw[arr] (tWgt) -- (tEqWgt);
\end{tikzpicture}
\end{minipage}\par\bigskip
\caption{Use of proved results. Parentheses in the group headings refer to
the enumeration in the abstract. Repeated nodes
(P\,\ref{p:KBpositive}, P\,\ref{p:KBshared}, P\,\ref{p:ColWie},
T\,\ref{t:ActF}, T\,\ref{t:KBhom}, T\,\ref{t:Net},
T\,\ref{t:KBoverlapCusp}) identify links between panels.
The reduction tools are consolidated in T\,\ref{t:Spe};
the quartic construction and minimum-degree result form the single
node P\,\ref{p:Qua}.
Arrows indicate proof usage only; panels are thematic, not connected components.
Nodes without incident arrows have independent proofs and no recorded downstream use.}
\label{fig:dep}
\end{figure}
\iffalse
\paragraph{Definitions and notation.}
Definitions~\ref{d:Siphons} and~\ref{d:Siphon} specify siphons;
Definition~\ref{d:pos} specifies positivity;
Definition~\ref{d:prMod} gives the production and threshold models;
Definition~\ref{d:Odd} gives odds responses;
Definition~\ref{d:ExZ} gives exact zero sets;
Definition~\ref{d:Inv} gives the invasion exponent;
Definition~\ref{d:FSW} gives the counting and stability weights;
Definition~\ref{d:LinRed} specifies invariant group reductions. They are prerequisites for understanding the corresponding
statements, but are not nodes in the proof-dependency graph.
\fi
\iffalse
\paragraph{Two multistability constructions.}
Theorems~\ref{t:KBsharp} and~\ref{t:siphMulti}  prescribe response values and vanishing derivatives to construct
arbitrarily many stable equilibria. The former fixes a slope bound and
finds the sharp uniform NGM bound; the latter fixes the response zero set,
with no common slope bound, preserving all siphons and their transverse
matrices. For its production--loss splitting, the boundary next-generation
matrices vanish.

\paragraph{Applications and repeated arguments.}
The persistence application after Theorem~\ref{t:KBexchange} combines that
theorem with Proposition~\ref{p:sadCus}; it is not a prerequisite of the
general proposition. Theorem~\ref{t:cuspAsym} identifies its two branches
with those of Theorem~\ref{t:KBexchange}, but proves its formulas and
nondegeneracy directly.

Theorems~\ref{t:siphMulti} and~\ref{t:KBoverlapCusp} compute their siphon
conditions or transverse matrices directly, rather than invoking
Proposition~\ref{p:InvF}.
Theorem~\ref{t:Spe}(2) uses Proposition~\ref{p:KBpositive} for
the extremal-equilibrium argument; its other parts supply spectral
preservation and bifurcation lifting.
The transverse application following Theorem~\ref{t:EqWgt}
uses Proposition~\ref{p:VarCert} and Theorem~\ref{t:Wgt}.


\paragraph{Standalone results.}
Proposition~\ref{p:KBrepresentation} extends the scope of the production
representation and is retained for that conclusion. The longer quotient,
diffusion and synchronization material is omitted from this short-paper
diagram; the former separate reduction and lifting results form one node.
\fi
\iffalse
\paragraph{Definitions and notation.}
Definitions~\ref{d:Siphons} and~\ref{d:Siphon} specify siphons;
Definition~\ref{d:pos} specifies positivity;
Definition~\ref{d:prMod} gives the production and threshold models;
Definition~\ref{d:Odd} gives odds responses;
Definition~\ref{d:ExZ} gives exact zero sets;
Definition~\ref{d:Inv} gives the invasion exponent;
Definition~\ref{d:FSW} gives the counting and stability weights;
Definition~\ref{d:LinRed} specifies invariant group reductions. They are prerequisites for understanding the corresponding
statements, but are not nodes in the proof-dependency graph.
\fi

